\chapter*{SGA 4、SGA $4\frac12$ 与 SGA 5 的阅读线索}
\addcontentsline{toc}{chapter}{SGA 4、SGA 4½ 与 SGA 5 的阅读线索}
\markboth{SGA 阅读线索}{SGA 阅读线索}
\origpage{7}{1}
SGA 4 的报告 I 至 VI 介绍\term{Grothendieck 拓扑}{Grothendieck topology}{topologie de Grothendieck}的一般理论，论述极为详细。在研究特殊拓扑，例如产生\term{晶体上同调}{crystalline cohomology}{cohomologie cristalline}的拓扑时，这些报告可能很有价值。平展拓扑与经典直觉十分接近，并不需要如此庞大的防护设施：只要熟悉例如 Godement 的书 [4]，再有一点信心，就够了。其他可用的参考资料是 Artin 讲义 [1] 的第一至三章、Giraud 的 Bourbaki 报告 [3]，或者 [Arcata] I。报告 VII、VIII 开始研究平展拓扑，内容比 Artin 的第三章和 [Arcata] II 更详细。

曲线的研究是平展上同调的关键。Artin 在 SGA 4 IX 中开启了这一研究，X 则给出有关\term{上同调维数}{cohomological dimension}{dimension cohomologique}的推论。要点在 [Arcata] III 中重述。报告 XI 对后续内容并非必需；它对光滑复代数簇这一特殊情形，以优美的方式证明了平展上同调与经典上同调之间的比较定理。

报告 XII、XIII 证明了固有态射的基本\term{基变换定理}{base change theorem}{théorème de changement de base}。正如 Artin 在 \textit{Algebraic approximation of structures over complete local rings}, Publ. Math. IHES 36 (1969), pp.~23--58, §3 中指出的，他的逼近定理可以简化这一证明。[Arcata] IV 采用了这条途径，但只概述证明中的多次分解归约，并极简略地指出报告 XIV 的应用。这个定理使我们能够定义“\term{紧支撑上同调}{cohomology with compact supports}{cohomologie à supports compacts}”和“\term{固有支撑高阶正像}{higher direct image with proper support}{image directe supérieure à supports propres}”函子 $R^if_!$。这些函子的性质归结为固有态射的高阶正像函子的性质。这是篇幅冗长的报告 XVII 的主题；要点已在 [Arcata] IV 中说明，不过 XVII 的某些部分也可能有独立用途。

\origpage{8}{2}
报告 XV、XVI 围绕光滑态射的基本\term{局部无环性定理}{local acyclicity theorem}{théorème d'acyclicité locale}展开。在相应但较简略的 [Arcata] V 中，可以看出 SGA 7 的影响；那里也给出了所需的 SGA 2 中结果的证明。

[Arcata] VI 专论\term{Poincaré 对偶}{Poincaré duality}{dualité de Poincaré}，比 SGA 4 XVIII 更清楚，但并不试图给出完整证明。对于坚持要求纯代数证明的读者，引用 [Dualité] §3 可以使论证简化；该篇不证明迹态射所满足的形式体系（XVIII §2）。

平展上同调中有一种与凝聚上同调相似的\term{对偶形式体系}{duality formalism}{formalisme de dualité}。Grothendieck 原来通过奇点消解和\term{纯性猜想}{purity conjecture}{conjecture de pureté}来建立它；猜想的陈述见 [Cycle] 2.1.4。SGA 4 XVI 在相对情形下证明了它，而 SGA 4 XIX 则在假定奇点消解的前提下处理了等特征情形。[Th. finitude] 用另一种方法，对维数为 $0$ 或 $1$ 的正则概形上的有限型概形，建立了其中的关键点。SGA 5 I 给出了若干进一步发展。SGA 5 III 说明，这套形式体系如何推出十分一般的 Lefschetz--Verdier 迹公式。

可见，在 SGA 5 的原始版本中，Lefschetz--Verdier 公式的建立仍依赖猜想，而且其中并没有计算局部项。对于 $L$ 函数的应用，这次讨论班在 Frobenius 态射这一特殊情形下给出了另一种完整证明，正是 [Rapport] 中的证明。关于定理的陈述和分解归约的轮廓，还可参阅 Grothendieck 的 Bourbaki 报告 [5]；关于 Grothendieck 如何将关键情形化到 Weil 已经处理的情形，简述见 [2] §10；关于后一情形的 $\ell$ 进处理，见 [Cycle] §3。

SGA 5 还详细论述了从 $\mathbb Z/\ell^n$ 系数上同调到 $\mathbb Z_\ell$ 系数上同调的过渡（J. P. Jouanolou 的报告 V、VI），以及 Frobenius 态射的形式体系（C. Houzel 的报告 XV）。

\origpage{9}{3}
由 I. Bucur 执笔的报告 X、XII 计算了曲线上层的 Euler--Poincaré 特征数，以及某些对应在其上同调上诱导的自身态射的迹。在报告 VII 中，Jouanolou 计算了一些经典概形的上同调，并给出应用。此外，Serre 的《Brauer 理论导论》已收入他的书 [6] 的第三部分。

\section*{参考文献}
\begin{description}
\item[{[1]}] M. Artin，\textit{Grothendieck topologies}，Harvard University，1962。
\item[{[2]}] P. Deligne，\textit{Les constantes des équations fonctionnelles des fonctions L}。载 \textit{Modular functions of one variable II}，pp.~501--597，Lecture Notes in Math.~349，Springer-Verlag。
\item[{[3]}] J. Giraud，\textit{Analysis situs}，Séminaire Bourbaki 256，1963 年 5 月，Benjamin。
\item[{[4]}] R. Godement，\textit{Topologie algébrique et théorie des faisceaux}，Act. Sci. et Ind.~1252，Hermann，1958。
\item[{[5]}] A. Grothendieck，\textit{Formule de Lefschetz et rationalité des fonctions L}，Séminaire Bourbaki 279，1964 年 12 月，Benjamin。
\item[{[6]}] J. P. Serre，\textit{Représentations linéaires des groupes finis}，Collection Méthodes，Hermann，1967。
\item[SGA] \textit{Séminaire de Géométrie algébrique du Bois-Marie}。
\item[SGA 4] \textit{Théorie des topos et cohomologie étale des schémas}，M. Artin、A. Grothendieck 和 J. L. Verdier 主持。Lecture Notes 269、270、305。
\item[SGA 5] \textit{Cohomologie $\ell$-adique et fonctions L}，A. Grothendieck 主持。即将收入 Lecture Notes 出版。
\end{description}
